\section{String Theory and Loop Quantum Gravity as HC Sectors}
\begin{quote}\small
\noindent\textbf{Status.} \emph{Legacy comparative-physics block inside the active main body.}\par
\noindent\textbf{Block function.} This section places two familiar high-energy frameworks into the compensator-based sector reading without promoting them to separate primitives.
\end{quote}
\label{sec:stringlqg}
%% ============================================================

\begin{theorem}[String--LQG duality via the Heisenberg compensator]
\label{thm:string-lqg-duality}
The Hardy space $\Hplus$ decomposes under the Heisenberg compensator
into two sectors:
\[
  \Hplus = \Hplusplus\oplus\Hplusminus,
\]
\[
  C(f) = \tfrac{1}{2}(f+Jf)\in\Hplusplus
  \quad\text{(symmetric, bosonic --- String theory)},
\]
\[
  A(f) = \tfrac{1}{2}(f-Jf)\in\Hplusminus
  \quad\text{(antisymmetric, fermionic --- LQG)}.
\]
Both sectors are unitarily equivalent by Stone--von Neumann
\cite{brendecke2026rh}: they are two representations of the same
Weyl algebra $\WA$, seen from different frames.
String theory and LQG are therefore not competing theories of quantum
gravity: they are the bosonic and fermionic projections of the same
object under the Heisenberg compensator.
\end{theorem}

\begin{proposition}[Shared vacuum and Planck floor]
\label{prop:shared-vacuum-planck}
Both sectors share the same vacuum and the same minimal scale:
\begin{itemize}
  \item \textbf{Vacuum:} $\zminus=0$ (the photon/Sphaera) is the ground
        state of both $\Hplusplus$ and $\Hplusminus$.
  \item \textbf{Planck floor:} $\lP=|\zminus|_{\min}$ is the string
        length in $\Hplusplus$ and the area quantum
        $\sqrt{A_{\min}}\sim\lP$ in $\Hplusminus$.
        Both emerge from the same threshold of Theorem~1.2(i).
\end{itemize}
String theory and LQG are therefore not competing theories of quantum
gravity: they are the bosonic and fermionic projections of the same
object under the Heisenberg compensator.
\end{proposition}

\begin{theorem}[Wick rotation as splitting mechanism]\label{thm:wick}
Under the imaginary proper time transformation $t\mapsto i\tau_E$
(Wick rotation), the biquaternion time unit transforms as:
\[
  it \;\mapsto\; i\cdot(i\tau_E) = -\tau_E \in\mathbb{R}_{<0}.
\]
This is precisely the third triolic sector $\gamma_3<0$: the Wick
rotation identifies the imaginary time component with the unstable
tachyon sector.
Consequently:
\begin{enumerate}[label=(\roman*)]
  \item The triolic balance condition
        $\gstring+\gamma_{\mathrm{LQG}}+\gamma_3=0$ is the stability
        condition under imaginary proper time: without it, the system
        diverges to $\pm\infty$.
  \item The two positive sectors ($\gstring,\gamma_{\mathrm{LQG}}>0$)
        are stable under Wick rotation; the negative sector
        ($\gamma_3<0$) is unstable and is projected away by the
        cascade $\KC$.
  \item The Minkowski norm $|q_n|^2=\tau_n^2$ becomes Euclidean
        $|q_E|^2=\tau_E^2\geq 0$ under $t\mapsto i\tau_E$: the system
        is unconditionally stable in Euclidean signature.
\end{enumerate}
\begin{proof}
(i). Under $t\mapsto i\tau_E$: $it\mapsto -\tau_E$.
The biquaternion $q_n=it\cdot\mathbf{1}+z_1 e_1+z_2 e_2+z_3 e_3$ becomes
$q_E=-\tau_E\cdot\mathbf{1}+z_1 e_1+z_2 e_2+z_3 e_3$
with norm $|q_E|^2=\tau_E^2+|z_1|^2+|z_2|^2+|z_3|^2\geq 0$.
Stability requires $\tau_E^2\geq 0$, which is always satisfied.
The balance $\gamma_1+\gamma_2+\gamma_3=0$ is equivalent to requiring
that the $-\tau_E$ component does not introduce a net positive energy
divergence.

(ii). In Euclidean signature, only $\gamma>0$ sectors contribute to a
convergent partition function $Z=\mathrm{Tr}(e^{-\beta H})$ with $H\geq 0$.
The sector $\gamma_3<0$ would give $e^{+\beta|\gamma_3|H}\to\infty$:
it is projected away by $\KC$.

(iii). Direct computation of the Euclidean norm.
\end{proof}
\end{theorem}

\begin{theorem}[Triolic structure of $\gBI$]
\label{thm:triolic-gBI}
The Barbero--Immirzi parameter $\gBI$ and the string-sector candidate
$\gstring = \ln 2/(\pi\sqrt{3})$ are not in contradiction.
They are two projections of the same triolic invariant:
\[
  \gstring + \gamma_{\mathrm{LQG}} + \gamma_3 = 0,
\]
with $\gamma_3 = -(\gstring+\gamma_{\mathrm{LQG}})\approx -0.3649$ being
the third (unoccupied, tachyon) sector.
The frame-invariant quantity is the triolic norm:
\[
  \tau_\gamma := \frac{1}{\sqrt{3}}
  \sqrt{\gstring^2+\gamma_{\mathrm{LQG}}^2+\gamma_3^2}\approx 0.2619.
\]
\end{theorem}

\begin{theorem}[Observer frame offset]\label{thm:frame-offset}
Every matter observer carries a nonzero contravariant component
$\zminus\neq 0$.
All measured physical constants --- including $\hbar$, $G$, $c$, and
$\gBI$ --- are projections of their frame-invariant values onto the
observer's sector:
\[
  \gamma_{\mathrm{observed}}
  = \gamma_0\cdot\frac{1}{\sqrt{1-\beta^2}}
  = \gamma_0\cdot\gamma_{\mathrm{Lorentz}}(\beta).
\]
The Barbero--Immirzi parameter $\gBI\approx 0.2375$ is not an absolute
constant of nature: it is the Lorentz-transformed $\gstring$, observed
from a matter frame with $\beta\approx 0.844$.
\end{theorem}

\begin{proof}
$\gstring = \ln 2/(\pi\sqrt{3})$ is computed in the photon frame
$\zminus=0$ from the HC trace formula.
The ratio $\gBI/\gstring = 1.864$ equals a Lorentz factor
$\gL = 1/\sqrt{1-\beta^2}$ for $\beta=\sqrt{1-1/1.864^2}\approx 0.844$.
The corresponding observer position is $\zminus=\beta/2\approx 0.422$.
\end{proof}

\begin{theorem}[Symmetric matter--tachyon frame]\label{thm:symmetric-frame}
Matter and tachyon are positioned symmetrically around the photon frame
$\zminus=0$:
\[
  \beta_{\mathrm{matter}} = +\bsym,\qquad
  \beta_{\mathrm{tachyon}} = -\bsym,
\]
where
\[
  \bsym = \sqrt{\frac{\gBI-\gstring}{\gBI+\gstring}}\approx 0.5493.
\]
The Lorentz factor between the matter and tachyon frames is:
\[
  \gL = \frac{1+\bsym^2}{1-\bsym^2}
      = \frac{\gBI}{\gstring}\approx 1.864.
\]
\end{theorem}

\begin{proof}
Let $r:=\gBI/\gstring$. Set $\bsym:=\sqrt{(r-1)/(r+1)}$. Then:
\[
  \frac{1+\bsym^2}{1-\bsym^2}
  =\frac{1+(r-1)/(r+1)}{1-(r-1)/(r+1)}
  =\frac{(r+1)+(r-1)}{(r+1)-(r-1)}=\frac{2r}{2}=r.\quad\checkmark
\]
The relativistic velocity of the tachyon as seen from the matter frame:
\[
  \beta_{\mathrm{tachyon}|\mathrm{matter}}
  =\frac{-\bsym-\bsym}{1-(-\bsym)\bsym}
  =\frac{-2\bsym}{1+\bsym^2},
\]
with corresponding Lorentz factor $\gamma_{\mathrm{tachyon}|\mathrm{matter}}=r$.\quad$\checkmark$
\end{proof}


%% ============================================================

